\section{State, token, and cycle commands}

\subsection{\code{SET} --- initialize or alias (\implemented)}

\begin{lstlisting}
SET Q 0p
SET Q 1p
SET Q sp
SET Local $Parameter
\end{lstlisting}
\code{SET target value} requires a positional value. A constant schedules or
emits initialization: zero, X from zero, or H from zero. A nonconstant value
makes the target an alias of that source. A \code{state} parameter's constant
SET is treated as a runtime default, allowing the UI or truth-table runner to
supply its actual start state.

\textbf{Syntax:} exactly one target and one value are meaningful. SET has no
\code{-I} or \code{-O} flags. It does not perform classical assignment at each
simulation step; it establishes compilation-time preparation or reference
identity. See the dedicated constants and aliases chapters before using SET on
unknown quantum data.

\subsection{\code{CREATETOKEN} --- reserve wires (\implemented)}

\begin{lstlisting}
CREATETOKEN -I Sum Carry Workspace
\end{lstlisting}
Every \code{-I} token immediately claims a stable wire index. No \code{-O} is
used. This is useful for declared outputs before their first gate.

\code{CREATETOKEN} does not initialize amplitudes by itself. Pair new workspace
with SET before relying on its value. Listing outputs early also makes process
structure easier to review:
\begin{lstlisting}
CREATETOKEN -I Sum Carry
SET Sum 0p
SET Carry 0p
\end{lstlisting}

\subsection{\code{DELETETOKEN} --- end a token name (\implemented)}

\begin{lstlisting}
DELETETOKEN -I Workspace
\end{lstlisting}
The named token leaves scope. A later use is a compile error. The wire is
returned to the allocator only when it is already known to be $\ket0$, which
is the state after \code{SET wire 0p} and before any gate targets it. Deletion
does not emit \code{RESET} and does not erase an unknown or entangled state.

\subsection{\code{FREE} --- same lifetime rule (\implemented)}

\begin{lstlisting}
FREE -I ScratchA ScratchB
\end{lstlisting}
FREE uses the same rule as DELETETOKEN. Use \code{SET wire 0p} when the
intended behavior is zero preparation. A wire that is not known $\ket0$ keeps
its qubit index, and only the name dies.

\subsection{\code{INCREASECYCLE} --- cycle boundary (\implemented)}

\begin{lstlisting}
INCREASECYCLE
\end{lstlisting}
The compiler flushes pending zero-initializations into an internal batched
\code{RESET}, increments the cycle counter, and prepares workspace for following
operations. Cycle suffixes on tokens do not create separate visual wires.

Cycles organize protocol stages. They are not wall-clock timing guarantees and
do not cause gates to execute concurrently. The final circuit is still an
ordered gate sequence. Each emitted gate records the timeline cycle, and the
canvas draws a dashed slice at \code{INCREASECYCLE}. An integer suffix such as
\code{Q:1} must match this process's cycle. A child process starts again at
cycle 0, so its own \code{:0} labels stay valid after the parent advances.
Non-numeric suffixes stay labels. The suffix does not allocate another wire.

\subsection{Cycle-based state evolution}
\hypertarget{cycle-based-state-evolution}{}

Cycle does \emph{not} mean loop. Cycle does not mean parallel execution. Cycle
does not mean wall-clock time. Cycle means the next logical stage of the same
finite circuit:

\begin{center}
\texttt{Cycle 0} \quad$\rightarrow$\quad \texttt{INCREASECYCLE}
\quad$\rightarrow$\quad \texttt{Cycle 1}
\quad$\rightarrow$\quad \texttt{INCREASECYCLE}
\quad$\rightarrow$\quad \texttt{Cycle 2}
\end{center}

A classical sequential rule such as $Q_{\mathrm{next}}=Q\oplus T$ is shown by
\emph{unrolling}:
\begin{lstlisting}
CNOT -I T -O Q
INCREASECYCLE
CNOT -I T -O Q
INCREASECYCLE
CNOT -I T -O Q
\end{lstlisting}
Conceptually $Q_0\to Q_1\to Q_2\to Q_3$. There is no backward execution edge.
The Interactive Workbench button \textbf{Add cycle boundary} inserts the same
marker visually. The starter circuit \emph{T state transition across cycles}
demonstrates this pattern.

\subsection{\code{JOIN} --- name an ordered register (\implemented)}

\begin{lstlisting}
JOIN -I A B -O AB
\end{lstlisting}
JOIN gives the ordered wires A then B the name AB. It is not a gate and
allocates no amplitudes: the simulator state is already one vector in
$(\mathbb{C}^2)^{\otimes n}$. At least two inputs and exactly one output are
required. A one-qubit gate on a multi-wire register is an error; name one wire
or \code{SPLIT} it first.

\subsection{\code{SPLIT} --- name a subsystem (\implemented)}

\begin{lstlisting}
SPLIT AB A 2
\end{lstlisting}
The third argument is a power-of-two dimension, the same convention as
\code{_dim}. \code{2} peels off one qubit from the front of AB and names it A.
The remaining wires stay on AB. SPLIT does not copy amplitudes. An entangled
register still has no separate pure state for A; the new name is a view of
those wires inside the global vector.

